\documentclass{article}

\usepackage[letterpaper]{geometry}
\usepackage[colorlinks=true, allcolors=blue]{hyperref}
\usepackage[spanish, activeacute]{babel}
\usepackage{tkz-euclide}
\usepackage{amsmath}
\usepackage{amsfonts}
\usepackage{amsthm}
\usepackage{amssymb,times}
\usepackage{graphicx}
\usepackage{latexsym}
\usepackage{subcaption}
\usepackage{xcolor}
\usepackage{pifont}
\usepackage{bbm}
\usepackage{pstricks}
\usepackage{pstricks-add}
\usepackage{pst-all}
\usepackage{mathrsfs}
\usepackage{pgfplots}
\usepackage{longtable}
\usepackage{multicol}
\usepackage{multirow}
\usepackage{kantlipsum}
\usepackage[inline]{enumitem}
\usepackage[export]{adjustbox}
\usepackage{array}
\usepackage{float}
\usepackage{booktabs}
\usepackage{mdframed}
\usepackage{minted}
\usepackage{tabularx}
\usepackage{adjustbox}
\usepackage[dvipsnames]{xcolor}


\usepackage{caption}
\captionsetup[sub]{labelformat=empty}

\definecolor{bgcolor}{RGB}{240,245,255}

\usetikzlibrary{angles, quotes, positioning}
\renewcommand\thesection{\arabic{section}}
\renewcommand{\proofname}{Demostración}
\renewcommand{\qedsymbol}{$\blacksquare$}
\newcommand{\pa}[1]{\left(#1\right)}
\newtheorem{lemma}{Lema}

\newcommand{\R}{\mathbb{R}}
\newcommand{\N}{\mathbb{N}}
\newcommand{\Q}{\mathbb{Q}}
\newcommand{\Z}{\mathbb{Z}}

\newenvironment{solution}
{\par\noindent\textit{Solución.}\ }
{\hfill$\blacktriangle$\par}

\begin{document}
    
    \begin{table*}
        \centering
        \begin{tabular}{cc}`'
            \adjustbox{valign=c}{\includegraphics[width=0.12\linewidth, valign=m]{logo-ug.png}} & 
            \adjustbox{valign=c}{
                \begin{tabular}{c}
                    {\large\uppercase{\textbf{Univers`~idad de Guanajuato}}}\\
                    {\small\uppercase{División de Ciencias Naturales y Exactas}}\\
                    {\small\uppercase{Campus Guanajuato}}
                \end{tabular}
            }
        \end{tabular}
    \end{table*}

    \begin{center}
        \textbf{Tarea 3 (Estructuras de Datos y Algoritmos I)}\\

        \begin{table*}
            \centering
            \begin{tabularx}{1\textwidth}{|X|X|X|}
                \hline
                \multicolumn{3}{|l|}{\textbf{Nombre:} Axel Magaña Falcón.}\\[2pt]
                \hline
                \textbf{Grupo: } 2$^\circ$A & \textbf{Fecha: }17/03/2026 & \textbf{Calificación:}\\[2pt]
                \hline
                \multicolumn{3}{|l|}{\textbf{Profesor:} Claudia Esteves.}\\[2pt]
                \hline
            \end{tabularx}
        \end{table*}
    \end{center}

    \begin{enumerate}[label=\textbf{Problema \arabic*}]
        \item \textbf{(Bubble Sort)} [\textcolor{ForestGreen}{\textbf{Accepted}}].
    
        Sobre este problema hay poco qué comentar. Es una implementación estándar de Bubble Sort e imprimir el arreglo tras cada intercambio de elementos.

        \item \textbf{(Insertion Sort)} [\textcolor{ForestGreen}{\textbf{Accepted}}].
        
        En este problema tampoco hay nada qué comentar. Es una implementación estándar de Insertion Sort, imprimiendo el arreglo tras cada cada paso de inserción. 

        \item \textbf{(Inversion Count)} [\textcolor{ForestGreen}{\textbf{Accepted}}].
        
        \textit{La solución estándar \textbf{(Merge Sort)}}: Para calcular las inversiones de un arreglo, es posible hacerlo mientras se ordena utilizando Merge Sort. Supongamos que nos encontramos en un paso intermedio a punto de hacer un merge de dos subarreglos. Supongamos que, durante el merge, agregamos un elemento del subarreglo izquierdo y hemos agregado previamente $k$ elementos del subarreglo derecho. Puesto que dichos $k$ elementos se encontraban a la derecha en el orden del arreglo original, nos permite contabilizar $k$ inversiones. Dicho proceso, en cada ejecución del proceso de merge nos permite contabilizar todas las inversiones en un tiempo de $\mathcal{O}(n \log n)$.

        \textit{Mi solución \textbf{(Árbol de Fenwick)}}: El Árbol de Fenwick es una estructura que nos permite calcular la suma de cada prefijo de un arreglo y actualizar un valor del mismo en un tiempo $\mathcal{O}(\log n)$. Puesto que el problema nos asegura un $A_i$ máximo de $10^7$, podemos utilizar un Árbol de Fenwick de tamaño $10^7$ para contabilizar los elementos existentes en el prefijo hasta $i$ del arreglo. Al iterar sobre el elemento $i + 1$ podemos calcular los elementos mayores a él que se encuentran detrás, esto nos permite encontrar las inversiones tales que $A_j > A_i$ para todo $0 \leq j \leq i$. Finalmente, actualizas la estructura contabilizando el elemento $i$ y dicho proceso se repite iterando sobre cada elemento, lo que nos permite contabilizar las inversiones en un tiempo de $\mathcal{O}(n \log n)$.
    \end{enumerate}
\end{document}